% a06_WN_d: part 4 of 4 of the walled reading of Ba3 Prop. 1 (input from sec/a06_BALtorus_d7b.tex): the proof of thm:a06:WN(ii), the constants, and two remarks on scope.

\begin{proof}[Proof of \cref{thm:a06:WN}\textup{(ii)}]
The proof of \cite[(1.33)]{Ba3} at $L=5$ for the bulk model, Steps~3(D$_3$)--3(D$_4$) with the twenty sketched steps listed after \cref{thm:a06d:S49}, is run for the walled model and for a class $\cC$ closed under the moves. It is organised as follows: first, where the scalar lines enter and what replaces each entry; second, where the class enters and what supplies each partner; third, the assembly and the induction.

\emph{Step 1: the scalar lines.} At zero background the scalar lines of every graph, slice, counterterm and cut kernel are zero-field kernels, and the proof uses them in exactly four ways. \emph{(1) Kernel bounds}: the group-(A) estimates (the slice bounds \cite[(2.6), (2.10)--(2.12)]{Ba3}, the kernel bounds of the cube step, the two-domain differences $\delta G_s(\square,\eta\Z^3,0)$, the cut-kernel estimate \cite[(3.1)]{Ba3} in the cluster form of \cref{cor:a06d:S24}, the Combes--Thomas bounds of the chains), each a composition of the pointwise bounds on the zero-field kernels with lattice sums; for walled lines these are \cref{lem:a06:WN1}(b),(c),(e), with constants polynomial in $\bar v$. The two-domain difference is bounded by the same resolvent identity as in the bulk, with \cref{lem:a06:WN1}(b) on each domain, since at zero field the bounds hold at every point of a Neumann domain and of the torus. \emph{(2) Algebra}: the slice identity, its rescaling, the gauge covariance used in the gauge-out of $B_0$, the cancellation of the operator in the differences of the $\tilde B'$-expansion and of the domain changes; for walled lines these are \cref{lem:a06:WN1}(a), and at zero background the $\tilde B'$-expansion, the constant field $B_0$ and its gauge-out are void. \emph{(3) The regrouped lines of the renormalisation}: in Step~3(D$_4$) the cumulative regrouped line of a node is, by the slice identity, the inverse of the operator pinned at the lowest external scale, which after rescaling to $\xi\Z^3$ is $\breve G_l=C^\xi+R_l$ (\cref{thm:a06d:S48}, \cref{lem:a06d:S47}); for walled lines the rescaled operator carries the potential $s^2v_{\mathrm{per}}(s\,\cdot)\in[0,\bar v]$ (\cref{lem:a06:WN1}(a),(f)), so $\breve G^v_l=C^\xi+R^v_l$ with the \emph{same} $C^\xi$ and with $R^v_l$ a remainder line of \cref{lem:a06d:S47}(b) with constant $C_R(1+\bar v)^3$ (\cref{lem:a06:WN1}(d)); hence every step of \cref{thm:a06d:S48} that concerns terms with at least one $R$-line (the one-loop $D=-1$ cancellation, which uses only the bounds on the $R$-line and the forward-difference structure of the $C^\xi$-line; the re-slicing of \cref{lem:a06d:S47}; the bound function estimate \cite[(2.16)]{Ba3}) runs verbatim, and every all-$C^\xi$ expression, with the inner counterterm members re-expanded with their own lines, is literally the same expression as in the bulk model \cite[pp.~437, 441]{Ba3}. \emph{(4) The counterterm values}: the model's own, by definition (the walled model). The torus-versus-$\xi\Z^3$ identification of \cite[pp.~418--419]{Ba3} is \cref{lem:a06:WN1}(f), with the image terms bounded, as in the bulk, by \cite[(1.33)]{Ba3} for the classes of lower order into which they decompose (the induction below). The sharp-indicator lemma \cref{lem:a06d:S46} uses only the decay of the internal lines, supplied by \cref{lem:a06:WN1}(b); the useless-counterterm lemma \cref{lem:a06d:S45} and the \cite[(2.4)]{Ba3} lemma \cref{lem:a06d:S44} are kernel-free given the bounds; \cref{lem:a06d:S41,lem:a06d:S42} are identities. Subgraphs with no scalar slot and an odd number of vector slots vanish at $B=0$ also with walled lines, which are scalar kernels times $\mathbf 1_N$ (\cref{lem:a06:WN1}(a)). A finite \cite[(1.7)]{Ba3} vertex with an $x$-dependent coefficient never sits inside a dangerous node (\cref{lem:a06d:S25}(b), \cite[Cor.~2.3]{Ba3}), so such coefficients enter only through their suprema, as the statement counts them.

\emph{Step 2: the class.} The class enters the bulk proof only at the dangerous nodes, as the sum over a node's partners with the rest of the graph fixed (\cref{lem:a06d:S41}(a): the counterterm pairing is an identity of class sums; \cref{thm:a06d:S48}: the local coefficient is the class sum at the node; \cref{thm:a06d:S28}), at the cut counterterms of partnered subgraphs (\cref{lem:a06d:S25}, which needs \cite[(1.33)]{Ba3} for the completed pieces of lower order), at the class-component cube (one label multiset), and in the regrouping after the cut (\cref{thm:a06d:S23}, kernel-free, (i)). Node by node, with walled lines: \emph{two-scalar-leg nodes} are paired with their counterterm member (P), the exact identities \cite[(3.19)]{Ba3} ($D=0$) and \cite[(3.11)]{Ba3} ($D=-1$) against the model's own counterterm with the same lines, the same localisation and (Agr); the $D=-1$ first moment has its all-$C^\xi$ part bounded by \cite[(3.16)]{Ba3} and its $R^v$-part by $|\partial R^v\partial^*|\le C(1+\bar v)^3(\xi+\rho)^{-1}e^{-\delta\rho}$, so $|c_\mu|\le C(1+\bar v)^3$; reducible chains by \cref{rem:a06:WNCH}. \emph{Two-vector-leg nodes}: the per-label two-vector-leg class at the node, which (W) supplies; the zeroth moment of its all-$C^\xi$ part vanishes by \cref{lem:a06d:S26}, applied as in \cref{thm:a06d:S48} to the full lines $C^\xi$ of the class sum (the lemma allows any $x$-dependent $|\phi|^2$ coefficients, so the inner counterterm and finite-vertex coefficients, and the potential itself, are admissible coefficients; \cref{lem:a06d:S26} is thus verbatim in the walled model), the first moment of its all-$C^\xi$ part by the reflection identity \cite[(3.30)]{Ba3} on $C^\xi$, which is model-free; the $R^v$-parts by the $R$-line steps of \cref{thm:a06d:S48}, with the zero-momentum remainder vertex kept. \emph{$\phi\phi A$ nodes}: in the bulk, \cref{thm:a06d:S48} bounds the zeroth moment by the reflection identity \cite[(3.37)]{Ba3} over the per-label full $\phi\phi A$ class; in a class closed under the moves this is replaced by (iii): at a bridgeless node the (X)-partner is in $\cC$ and (iii)(a) cancels the pair exactly at a common anchor, or, when the node carries the pinned root at one of its two slot vertices, (iii)(b) bounds the pair by $C^v_7$; at a bridged node no partner is needed, (iii)(c), the counterterm member (P) of $S_2$ being in $\cC$; and the pair bounds are of the form $|c^{(0)}|\le C$ that \cref{thm:a06d:S49} consumes for a $\phi\phi A$ local vertex. \emph{\cite[(2.4)]{Ba3} subgraphs}: summation by parts, \cref{lem:a06d:S44}, kernel-free. \emph{Vacuum components}: cancelled by the model's own $E_s[v]$ with the same labels and lines, up to image terms (Step~1). \emph{One- and three-vector-leg subgraphs}: vanish (Step~1). The cut counterterms of partnered subgraphs (\cref{lem:a06d:S25}) are classes of lower order of the completed pieces, which are closed under the moves (a move at a divergent subgraph of a piece is a move of the uncut member, and the completion of a partnered subgraph is by the moves); and the class-component cube depends on the label multiset only, common to $\cC$. Per-label full classes are closed under the moves, so the bulk statement for them is the special case $\cC=$ full class; and every partner that the walled node mechanisms use is in $\cC$ by the closure, while every other step is a bound per term.

\emph{Step 3: assembly and induction.} With Steps~1 and 2, \cref{thm:a06d:S48,thm:a06d:S49} hold for $\cC$ with walled lines, with the $\phi\phi A$ clause in the pair form of (iii); the induction on the number of lines inside \cref{thm:a06d:S48} and on the number of vertices in \cref{lem:a06d:S25} runs over classes closed under the moves of lower order (the completed pieces of a cut partnered subgraph, and the pieces into which an image term decomposes), to which the statement applies by the induction hypothesis; the base is the one-vertex graphs. The bound function estimate \cite[(2.16)]{Ba3} then gives \cite[(3.3)]{Ba3} for every class component, and the assembly of \cite[pp.~420--445]{Ba3} listed after \cref{thm:a06d:S49}, with the cluster cut and the cube steps, gives \cite[(1.33)]{Ba3} in the stated form, with $d_v$, $d_s$ counting the action couplings, the counterterm vertices as exact counterterms, and the finite vertices through their suprema (they have $D(v)=2$ whatever their coefficient, \cref{lem:a06d:S25}(b)); the localisation of an external field whose leg moves under the moves is (H$_A$); the slot fields of several variables are measured by \cite[(1.32)]{Ba3} in the slot variable, as the two-variable external fields of \cite[p.~419]{Ba3}; and in a two-point setting the root vertex is fixed at $x$ throughout, every kernel bound used being pointwise, so the bound is uniform in $x$.

\emph{Constants.} The walled inputs carry $(1+\bar v)^2$ (\cref{lem:a06:WN1}(b)), $(1+\bar v)^3$ (\cref{lem:a06:WN1}(d)) and the constants of \cref{lem:a06:WN1}(e), polynomial in $\bar v$; a graph of order $\le\bar n$ has a number of lines bounded by $\bar n$, so $O(1)^v$ is a polynomial in $\bar v$ of degree bounded in terms of $\bar n$, and $\delta^v_0$ is a fixed fraction of $\delta^w$, $\delta_R$ and the bulk rates. The parameter rows of Step~3(D) are kept, with $R_0$ replaced by $R^v_0\defeq R_0+10$ in every depth row for a scalar propagator bound and the torus side $\ge\max\{6\sqrt3\,r(t_s)+2,8/\delta^v_0\}$, all lower bounds on $R$, $r^\flat$ and the torus side or caps on $\eps_0$ of Ba{\l}aban's kind, fixed after $R_0$ and before $\eps_0$; nothing in the order of the choices changes. The constants are therefore independent of $\eps$, of $\eps_0\in(0,\eps^*_0]$ below the caps, of $s$, $K$, the labels, the torus side, $x$, and of $v$ beyond $\bar v$, and they depend on $L=5$, on the parameters of Step~3(D) and polynomially on $\bar v$.
\end{proof}

\begin{remark}[non-circularity of the uses in the rotation part]\label{rem:a06:WNcirc}
\Cref{thm:a06:WN}(ii) is a bound in which the suprema of the finite-vertex coefficients are parameters. Its uses in the rotation part at coupling order $4$ carry, as finite vertices, only the non-renormalising parts of the order-$2$ inner counterterms, whose sizes come from the order-$2$ statement, and the wall vertex, whose coefficient is $\le nt_s^2$; the sizes of the order-$4$ finite vertices are then \emph{consequences} of (ii) at order $4$, not inputs to it. The induction of the proof runs on the number of vertices inside the theorem.
\end{remark}

\begin{remark}[a bounded wall is not a boundary]\label{rem:a06:WNbdry}
\Cref{thm:a06:WN} is stated for a bounded site potential on the scalar lines with the vector lines unwalled, on a torus. It does not cover scalar or vector lines with Dirichlet or Neumann boundary conditions on a box, and the proof of (iii)(b) does not extend to them: a boundary changes the kernel by image terms $C^\xi(x-y^*)$ which, for $x$ and $y$ in the boundary layer, are as singular as $C^\xi$ itself and so are not remainder lines in the sense of \eqref{eq:a06:WNR}; with a Neumann vector line the odd part of the pinned pair is left uncancelled at every scale, and nothing bounds $|c_{S,1}-c_{S,2}|$ uniformly in $\xi$ in the layer. A bounded potential changes the kernel by $R^v$, which satisfies \eqref{eq:a06:WNR} at \emph{every} pair of points (\cref{lem:a06:WN1}(d); the multiplier $1-w$ is never differentiated, and no smoothness of $v$ is used), so there is no boundary layer, and the vector line is the translation-invariant torus propagator. This is why the rotation part compares King's cubes with walled tori and never applies the renormalisation theorem on a cube.
\end{remark}
